\documentclass[12pt,a4paper]{article}

\usepackage[T2A]{fontenc}
\usepackage[utf8]{inputenc}
\usepackage[russian]{babel}
\usepackage{amsmath,amssymb,amsthm}
\usepackage{geometry}
\usepackage{hyperref}
\usepackage{bookmark}
\usepackage{mathtools}
\usepackage{enumitem}
\usepackage{microtype}

\geometry{margin=22mm}
\hypersetup{
  colorlinks=true,
  linkcolor=blue,
  urlcolor=blue,
  citecolor=blue
}
\setlist{itemsep=0.25em, topsep=0.35em}

\newcommand{\R}{\mathbb{R}}
\newcommand{\grad}{\nabla}
\newcommand{\norm}[1]{\left\lVert #1\right\rVert}

\title{Методы оптимизации --- подробный конспект к зачёту}
\author{}
\date{\today}

\begin{document}
\maketitle

\section*{Для кого этот текст и как им пользоваться}
Документ рассчитан на \textbf{первое изучение}: к каждому методу добавлены пояснения ``с нуля'', численные идеи и формулы, которые обычно ожидают на зачёте. В конце разделов иногда даются ссылки на Wikipedia как дополнительное чтение.

Каждый метод ниже можно читать по шаблону:
\begin{enumerate}[label=\textbf{\arabic*.}]
  \item \textbf{Зачем} --- какую задачу решаем и что на выходе.
  \item \textbf{Что нужно заранее} --- какие данные и предположения (гладкость, унимодальность и т.д.).
  \item \textbf{Идея простыми словами} --- без формул или почти без них.
  \item \textbf{Шаги алгоритма} --- что делать на одной итерации.
  \item \textbf{Формулы} --- то, что обычно спрашивают на зачёте.
  \item \textbf{Остановка и практика} --- когда заканчивать и на что смотреть в реализации.
\end{enumerate}
Если что-то кажется абстрактным, прочитайте сначала только пункты 1--3, затем вернитесь к формулам.

\tableofcontents
\bigskip

%--------------------------------------------------------------------
\section{Постановка задач оптимизации}
%--------------------------------------------------------------------

\subsection{Что такое оптимизация в одном абзаце}
Нужно выбрать \emph{решение} \(x\) (набор чисел: параметры модели, режим работы, маршрут и т.д.) так, чтобы некоторая \emph{целевая функция} \(f(x)\) была как можно меньше (или больше). Множество допустимых \(x\) может быть ограничено равенствами и неравенствами --- тогда говорят об \textbf{условной} оптимизации.

\subsection{Безусловная оптимизация}
\textbf{Задача:} найти точку (или набор точек), где \(f\) минимальна на всём \(\R^n\):
\[
  x^\star \in \arg\min_{x \in \R^n} f(x).
\]
\textbf{Важно:} минимум может не достигаться (например, \(f(x)=e^{-x}\) на \(\R\)). Тогда ищут \(\inf f\) и последовательность \(x_k\), для которой \(f(x_k)\) стремится к нижней грани.

\subsection{Условная оптимизация}
\textbf{Задача:}
\[
  x^\star \in \arg\min_{x \in \Omega} f(x),\qquad
  \Omega=\{x:\ g_i(x)\le 0,\ i=1,\ldots,m;\ h_j(x)=0,\ j=1,\ldots,p\}.
\]
\textbf{Смысл:} часть решений ``отбрасывается'' ограничениями; оптимум может лежать \emph{внутри} области или \emph{на границе} (тогда активны некоторые \(g_i=0\)).

\subsection{Локальный и глобальный минимум}
\begin{itemize}
  \item \textbf{Глобальный минимум} \(x^\star\): \(f(x^\star)\le f(x)\) для \emph{всех} допустимых \(x\). Это ``самая низкая точка на всей карте рельефа''.
  \item \textbf{Локальный минимум}: существует ``окрестность'' (маленький шар вокруг точки), внутри которой \(x^\star\) лучше всех. Вне окрестности могут быть точки с меньшим значением \(f\).
\end{itemize}
\textbf{Почему это важно:} большинство классических численных методов (градиентный спуск, Ньютон, BFGS) \emph{в лучшем случае} находят локальный минимум или стационарную точку. Чтобы искать глобальный минимум, подключают мультистарт, штрафы к глобальным эвристикам (PSO, ГА) или имитацию отжига.

%--------------------------------------------------------------------
\section{Теория: условия оптимальности и выпуклость}
%--------------------------------------------------------------------

\subsection{Градиент и матрица Гессе (напоминание)}
Если \(f\) гладкая, то \textbf{градиент} \(\grad f(x)\) --- это вектор частных производных: он указывает направление \emph{наибольшего} роста \(f\) в окрестности точки. \textbf{Гессиан} \(H(x)=\nabla^2 f(x)\) --- матрица вторых производных; он описывает ``кривизну'' функции в разных направлениях.

\subsection{Необходимые условия экстремума (без ограничений)}
Если \(x^\star\) --- локальный минимум и \(f\) дифференцируема в \(x^\star\), то
\[
  \grad f(x^\star)=0.
\]
\textbf{Интуиция:} в минимуме ``наклон'' в любом направлении не может быть отрицательным и положительным одновременно; значит линейная часть приращения \(f\) обнуляется.

Если \(f\) дважды дифференцируема, то дополнительно \(H(x^\star)\) \textbf{положительно полуопределена}: \(H(x^\star)\succeq 0\) (все собственные значения \(\ge 0\)).

\subsection{Достаточные условия (без ограничений)}
Если \(\grad f(x^\star)=0\) и \(H(x^\star)\succ 0\) (все собственные значения \(>0\)), то \(x^\star\) --- \textbf{строгий} локальный минимум. Интуитивно: функция ведёт себя как ``вогнутая параболоид вверх'' около точки.

\subsection{Выпуклая оптимизация}
Функция \(f\) \textbf{выпукла} на выпуклом множестве \(X\), если отрезок между любыми двумя точками графика лежит не ниже хорды:
\[
  f(\lambda x+(1-\lambda)y)\le \lambda f(x)+(1-\lambda)f(y),\quad \lambda\in[0,1].
\]
Для дифференцируемой выпуклой функции верно неравенство:
\[
  f(y)\ge f(x)+\langle \grad f(x), y-x\rangle.
\]
\textbf{Главный вывод для зачёта:} если \(f\) выпукла на \(X\), то любая точка с \(\grad f=0\) --- \textbf{глобальный} минимум на \(X\). Это сильно упрощает анализ.

\subsection{Условия ККТ (Каруша--Куна--Таккера)}
Рассмотрим
\[
  \min f(x)\quad \text{при}\quad g_i(x)\le 0,\quad h_j(x)=0.
\]
Лагранжиан:
\[
  L(x,\lambda,\nu)= f(x)+\sum_{i=1}^m \lambda_i g_i(x)+\sum_{j=1}^p \nu_j h_j(x).
\]
При стандартных условиях регулярности (например, LICQ: градиенты активных ограничений линейно независимы) в локальном минимуме \(x^\star\) существуют множители \(\lambda^\star,\nu^\star\) такие, что
\begin{align*}
  \grad_x L(x^\star,\lambda^\star,\nu^\star) &= 0,\\
  g_i(x^\star)&\le 0,\quad h_j(x^\star)=0,\\
  \lambda_i^\star &\ge 0,\\
  \lambda_i^\star g_i(x^\star) &= 0\quad\text{(дополняющая нежёсткость)}.
\end{align*}
\textbf{Смысл дополняющей нежёсткости:} если ограничение не активно (\(g_i(x^\star)<0\)), то \(\lambda_i^\star=0\); если активно (\(g_i=0\)), множитель \(\lambda_i^\star\) может быть положительным --- он ``подталкивает'' к допустимости.

%--------------------------------------------------------------------
\section[Одномерная оптимизация]{Одномерная оптимизация: поиск минимума на отрезке}
%--------------------------------------------------------------------
\textbf{Постановка:} \(\min_{x\in[a,b]} f(x)\).

\subsection{Унимодальность}
Функция \(f\) \textbf{унимодальна} на \([a,b]\), если существует точка \(x^\star\in[a,b]\) такая, что \(f\) строго убывает слева от \(x^\star\) и строго возрастает справа (в широком смысле --- не возрастает/не убывает). Тогда на \([a,b]\) один ``главный'' минимум; методы сужения отрезка опираются именно на это свойство.

\subsection{Метод перебора (сетка)}
\textbf{Зачем:} самый простой способ получить приближённый минимум, если не хочется думать про свойства \(f\).

\textbf{Идея:} разбить \([a,b]\) на равные отрезки шага \(h\), вычислить \(f\) в узлах, взять узел с наименьшим значением.

\textbf{Формулы:} узлы \(x_j=a+jh\), \(j=0,1,\ldots,N\), где \(N\approx (b-a)/h\).

\textbf{Оценка:} \(\Theta(1/h)\) вычислений \(f\). Точность по аргументу порядка \(h\).

\textbf{Минусы:} при малом \(h\) очень дорого; при большом \(h\) можно ``промахнуться'' мимо узкого минимума.

\subsection{Метод деления пополам (дихотомия)}
\textbf{Зачем:} сузить отрезок, содержащий минимум унимодальной функции, за логаритмическое число шагов.

\textbf{Идея:} возьмем две точки \(x_1<x_2\) внутри текущего отрезка \([a_k,b_k]\). Сравним \(f(x_1)\) и \(f(x_2)\). Унимодальность гарантирует: минимум \emph{не может} оказаться сразу с обеих сторон ``плохой'' половины --- отрезок можно сократить.

\textbf{Шаги:}
\begin{enumerate}
  \item Положить \(x_1=\frac{a_k+b_k}{2}-\delta\), \(x_2=\frac{a_k+b_k}{2}+\delta\) с малым \(\delta>0\).
  \item Если \(f(x_1)\le f(x_2)\), то новый отрезок \([a_k,x_2]\); иначе \([x_1,b_k]\).
  \item Повторять, пока длина отрезка \(>\varepsilon\).
\end{enumerate}
\textbf{Про \(\delta\):} без него \(x_1=x_2\) и сравнение бессмысленно; \(\delta\) берут много меньше требуемой точности по \(x\).

\subsection{Метод золотого сечения}
\textbf{Зачем:} как дихотомия, но на каждом шаге нужно только \textbf{одно новое} вычисление \(f\) (вторая точка ``наследуется'' с прошлого шага).

\textbf{Идея:} разместить две внутренние точки симметрично относительно концов отрезка так, чтобы при отбрасывании части отрезка одна из точек осталась внутри нового отрезка в том же ``золотом'' отношении.

\textbf{Формулы:} \(\tau=\frac{\sqrt5-1}{2}\approx 0{,}618\). Для \([a,b]\):
\[
  x_1=b-\tau(b-a),\qquad x_2=a+\tau(b-a).
\]
Если \(f(x_1)\le f(x_2)\), новый отрезок \([a,x_2]\); иначе \([x_1,b]\). Длина отрезка уменьшается примерно в \(\tau\) раз за итерацию.

\textbf{Сколько итераций до точности \(\varepsilon\).} Если начальная длина \(L_0=b-a\), то после \(k\) шагов \(L_k\approx \tau^k L_0\). Отсюда \(k\gtrsim \dfrac{\ln(L_0/\varepsilon)}{\ln(1/\tau)}\).

\textbf{Почему ``золотое''.} Отношение \(\tau\) --- единственное положительное число, для которого при симметричном размещении двух точек внутри отрезка одна из них автоматически попадает внутрь нового отрезка с тем же относительным положением --- поэтому на шаге нужна только одна новая оценка \(f\).

\subsection{Метод Фибоначчи}
\textbf{Отличие от золотого сечения.} Если заранее зафиксировано ровно \(N\) вычислений \(f\), метод Фибоначчи подбирает положения точек так, чтобы \textbf{минимизировать финальную длину интервала} после \(N\) измерений (для унимодальной функции).

\textbf{Числа Фибоначчи.} \(F_0=F_1=1\), далее \(F_{m+1}=F_m+F_{m-1}\).

\textbf{Идея размещения точек.} На первом шаге на отрезке длины \(L\) ставят две точки на расстояниях \(F_{N-1}/F_N\cdot L\) и \(F_{N-2}/F_N\cdot L\) от одного конца (формулы в разных учебниках записывают через отношения \(F_{N-k}/F_N\)). После сравнения \(f\) интервал сокращается, а \textbf{одна} из точек остаётся внутри нового интервала в нужном отношении для оставшегося числа измерений \(N-1\).

\textbf{Связь с золотым сечением.} При \(N\to\infty\) отношения \(F_{N-1}/F_N\) стремятся к \(\tau\), поэтому для ``очень большого числа шагов'' метод Фибоначчи близок к золотому сечению, но на конечном \(N\) он \emph{оптимален} по длине финального интервала.

\subsection{Метод Свенна (расширение интервала)}
\textbf{Зачем:} когда нет готового \([a,b]\), но есть стартовая точка \(x_0\) ``где-то рядом'' с долиной минимума.

\textbf{Идея:} идти из \(x_0\) в сторону убывания \(f\), увеличивая шаг (например, геометрически: \(h,2h,4h,\ldots\)), пока не получится тройка точек
\[
  x_{k-1}<x_k<x_{k+1},\qquad f(x_{k-1})\ge f(x_k)\le f(x_{k+1}),
\]
обрамляющая минимум. После этого к обрамляющему отрезку применяют золотое сечение.

\subsection{Параболическая интерполяция (квадратичная аппроксимация)}
\textbf{Идея:} по трём точкам с разными значениями \(f\) построить параболу \(q(x)\), совпадающую с \(f\) в этих точках, и перейти к её минимуму. Затем обновить тройку точек, сохраняя обрамление.

\textbf{Когда осторожно:} если \(f\) шумная или тройка плохая, парабола может вести куда угодно; в библиотеках обычно комбинируют с золотым сечением (метод Брента).

%--------------------------------------------------------------------
\section{\texorpdfstring{Многомерная безусловная оптимизация: методы спуска в \(\R^n\)}{Многомерная безусловная оптимизация: методы спуска}}
%--------------------------------------------------------------------
Общая схема итерации:
\[
  x_{k+1}=x_k+\alpha_k p_k,
\]
где \(p_k\) --- направление (часто ``вниз'' по рельефу \(f\)), \(\alpha_k>0\) --- длина шага.

\subsection{Направление спуска}
Вектор \(p_k\) называют \textbf{направлением спуска} в точке \(x_k\), если
\[
  \grad f(x_k)^\top p_k<0.
\]
Тогда при достаточно малом \(\alpha>0\) выполняется \(f(x_k+\alpha p_k)<f(x_k)\) (из разложения Тейлора первого порядка).

\subsection{Линейный поиск (подбор шага \(\alpha\))}
\textbf{Зачем.} Имеем направление \(p_k\) (например, \(-\grad f(x_k)\) или квазиньютоновское). Хотим \(\alpha>0\) такое, что \(f(x_k+\alpha p_k)\) существенно меньше \(f(x_k)\), но без ``прыжка'' через долину, где линейная модель уже лжёт.

\textbf{Модель первого порядка.} Локально \(f(x_k+\alpha p_k)\approx f(x_k)+\alpha\,\grad f(x_k)^\top p_k\). Если \(p_k\) --- направление спуска, то \(\grad f(x_k)^\top p_k<0\), и при малых \(\alpha\) значение убывает линейно по \(\alpha\).

\textbf{Условие Армихо} (достаточное убывание): для \(c_1\in(0,1)\) (часто \(10^{-4}\)) ищут \(\alpha>0\):
\[
  f(x_k+\alpha p_k)\le f(x_k)+c_1\,\alpha\,\grad f(x_k)^\top p_k.
\]
Правая часть --- это прямая ``ниже'' линейной модели; требование: реальное \(f\) должно оказаться под этой прямой.

\textbf{Backtracking (затухающий шаг).} Берут \(\alpha=\alpha_0\) (например, \(1\)), пока не выполнится Армихо, делают \(\alpha\leftarrow\rho\alpha\) с \(\rho\in(0,1)\) (например, \(0{,}5\)). Это дёшево: несколько вычислений \(f\).

\textbf{Сильное правило Вольфе (кривизна).} Часто добавляют второе условие с \(c_2\in(c_1,1)\) (часто \(0{,}9\)):
\[
  \bigl|\grad f(x_k+\alpha p_k)^\top p_k\bigr|\le c_2\,\bigl|\grad f(x_k)^\top p_k\bigr|.
\]
Упрощённый вариант (слабое Вольфе): \(\grad f(x_k+\alpha p_k)^\top p_k \ge c_2\,\grad f(x_k)^\top p_k\).

\textbf{Смысл.} Если производная \(\frac{d}{d\alpha}f(x_k+\alpha p_k)=\grad f(x_k+\alpha p_k)^\top p_k\) слишком отрицательна в конце шага, значит \emph{вдоль} \(p_k\) ещё есть крутой спуск --- шаг можно было продолжить, но часто это признак, что мы уже близко к ``дну'' и квазиньютоновская модель требует обновления; Вольфе балансирует длину шага.

\subsection{Градиентный спуск}
\textbf{Зачем:} базовый метод для гладких функций, когда есть градиент.

\textbf{Идея ``с нуля'':} в точке \(x_k\) функцию \emph{подменяют} самой простой нелинейной моделью --- \textbf{линейной} (это и есть разложение Тейлора \textbf{первого} порядка):
\[
  f(x_k+s)\approx f(x_k)+\grad f(x_k)^\top s.
\]
Если мы хотим \emph{сразу} сильно уменьшить эту линейную модель, логично идти в сторону, противоположную градиенту: \(s=-\grad f(x_k)\). Но линейная модель при уходе далеко от \(x_k\) перестаёт похожа на настоящую \(f\), поэтому делают \textbf{короткий} шаг \(\alpha_k\) и повторяют заново в новой точке.

\textbf{Формулы:} \(x_{k+1}=x_k-\alpha_k\grad f(x_k)\).

\textbf{Про постоянный шаг:} если \(\grad f\) липшицев с константой \(L\) (часто так для гладких функций с ограниченной второй производной), то при \(\alpha\in(0,2/L)\) метод с постоянным \(\alpha\) монотонно уменьшает \(f\) (в выпуклом случае --- классическая оценка).

\textbf{Проблемы:} в ``узких долинах'' траектории сильно осциллируют; сходимость может быть медленной.

\textbf{Как выбрать \(\alpha\) на практике.} Константный шаг (если известна \(L\)-гладкость), либо Армихо/backtracking на каждом шаге. Для выпуклых задач иногда используют Barzilai--Borwein или адаптивные правила --- это уже улучшения базовой схемы.

\subsection{Метод наискорейшего спуска}
\textbf{Идея.} Направление то же: \(p_k=-\grad f(x_k)\), но вместо ``угадывания'' \(\alpha_k\) решают одномерную задачу
\[
  \alpha_k=\arg\min_{\alpha\ge 0} \varphi(\alpha),\qquad \varphi(\alpha):=f\bigl(x_k-\alpha\grad f(x_k)\bigr).
\]
Функция \(\varphi(\alpha)\) --- это \(f\) на луче из \(x_k\) вдоль антиградиента. Её минимум --- \textbf{точка касания} уровня \(f\) лучом (в гладком выпуклом случае).

\textbf{Как искать \(\alpha_k\).} Одномерный метод (золотое сечение, параболы + Брент) или приближённый поиск с Вольфе.

\textbf{Плюсы.} За один ``большой'' шаг можно сильно уменьшить \(f\), если овраг не слишком кривой.

\textbf{Минусы.} Каждый шаг дороже; для плохо обусловленных квадратичных функций траектории всё равно могут зигзагообразно идти к минимуму --- это не лечится точным \(\alpha_k\) вдоль градиента.

\subsection{Метод сопряжённых градиентов (для квадратичной функции)}
Рассмотрим \(f(x)=\tfrac12 x^\top Ax-b^\top x\) с симметричной положительно определённой матрицей \(A\succ 0\).

\textbf{Идея ``с нуля''.} Градиент в точке \(x\): \(g(x)=Ax-b\). Минимум --- решение \(Ax=b\). Можно было бы решать СЛАУ напрямую, но если \(A\) большая и разреженная, полезны итерационные методы. CG строит направления \(p_0,p_1,\ldots\), которые \textbf{\(A\)-ортогональны} (сопряжены): \(p_i^\top A p_j=0\) при \(i\neq j\). Тогда каждый шаг уменьшает ошибку в ``новом'' направлении, не портя уже достигнутое вдоль предыдущих; для квадратичной функции за \(\le n\) шагов (точная арифметика) получается точное решение.

\textbf{Величины.} Невязка \(r_k:=g(x_k)=Ax_k-b\). Начало: \(p_0=-r_0\).

\textbf{Итерация (Флетчер--Ривс):}
\[
  r_k=Ax_k-b,\quad p_0=-r_0,
\]
\[
  \alpha_k=\frac{r_k^\top r_k}{p_k^\top A p_k},\quad x_{k+1}=x_k+\alpha_k p_k,
\]
\[
  r_{k+1}=r_k+\alpha_k A p_k,\quad
  \beta_{k+1}=\frac{r_{k+1}^\top r_{k+1}}{r_k^\top r_k},\quad
  p_{k+1}=-r_{k+1}+\beta_{k+1}p_k.
\]
\textbf{Замечание про \(\beta_k\).} В формуле Флетчера--Ривса \(\beta_{k+1}=\|r_{k+1}\|^2/\|r_k\|^2\). Другой популярный вариант --- Полак--Рибьер:
\[
  \beta_{k+1}^{\mathrm{PR}}=\frac{r_{k+1}^\top(r_{k+1}-r_k)}{\|r_k\|^2}.
\]
Для \textbf{нелинейного} \(f(x)\) направления перестают быть точно \(A\)-сопряжёнными; тогда на каждом шаге берут \(p_k\) по той же схеме, но \(r_k=\grad f(x_k)\), а \(\alpha_k\) находят линейным поиском (часто Вольфе), иногда с дополнительным условием, чтобы \(p_k\) оставался спуском.

\subsection{Метод Ньютона (подробно для начинающих)}
\textbf{Что мы вообще хотим.} В безусловной минимизации ищут точку, где градиент обнуляется: \(\grad f(x^\star)=0\). Это \textbf{не то же самое}, что ``найти корень произвольной функции'', но очень похоже: если ввести \(g(x)=\grad f(x)\), то для минимума нужно решить уравнение \(g(x)=0\). Классический метод Ньютона--Рафсона для уравнения использует производную; в оптимизации на роль ``производной'' выступает \textbf{матрица вторых производных} --- Гессиан.

\subsubsection{Одна переменная: от ряда Тейлора до формулы ``минус первая вторая''}
Зафиксируем текущую точку \(x_k\) и обозначим малое смещение \(t\) (шаг вдоль оси \(x\)). Разложим \(f\) в ряд Тейлора \textbf{до второго порядка} около \(x_k\):
\[
  f(x_k+t)\approx q_k(t):= f(x_k)+f'(x_k)\,t+\tfrac12 f''(x_k)\,t^2.
\]
Геометрически \(q_k(t)\) --- это \textbf{парабола} по \(t\). Если \(f''(x_k)>0\), парабола ``улыбкой вверх'' и у неё есть единственный минимум.

Чтобы найти минимум \(q_k(t)\), возьмём производную по \(t\) и приравняем нулю:
\[
  \frac{d}{dt}q_k(t)= f'(x_k)+f''(x_k)\,t = 0
  \quad\Rightarrow\quad
  t_k=-\frac{f'(x_k)}{f''(x_k)}.
\]
\textbf{Важная мелочь про ``первую делить на вторую'':} в учебниках часто говорят коротко ``шаг пропорционален производным'', но знак уже в формуле: это именно \(-f'/f''\), а не \(+f'/f''\). Минус означает: если \(f'(x_k)>0\) (функция локально растёт), нужно сдвинуться \emph{влево}, если \(f'<0\) --- \emph{вправо}.

Ньютоновская итерация в одномерии записывается так:
\[
  x_{k+1}=x_k+t_k=x_k-\frac{f'(x_k)}{f''(x_k)}.
\]
\textbf{Мини-пример.} Пусть \(f(x)=\tfrac12 a x^2+bx+c\) при \(a>0\). Тогда \(f'(x)=ax+b\), \(f''(x)=a\). Один шаг Ньютона из любой точки \(x_k\) даёт
\[
  x_{k+1}=x_k-\frac{ax_k+b}{a}=-\frac{b}{a},
\]
то есть \textbf{сразу} точный минимум параболы. Для общей гладкой функции парабола лишь \emph{приближает} \(f\) около \(x_k\), поэтому шаги повторяют.

\subsubsection{Много переменных: то же самое, но ``вторая производная'' --- матрица}
Пусть \(f:\R^n\to\R\) дважды дифференцируема. Вектор неизвестного шага обозначим \(s\in\R^n\). Разложение Тейлора второго порядка около \(x_k\):
\[
  f(x_k+s)\approx m_k(s):= f(x_k)+\grad f(x_k)^\top s+\tfrac12 s^\top H(x_k)s,
\]
где \(H(x_k)=\nabla^2 f(x_k)\) --- симметричная матрица (Гессиан) из вторых частных производных.

Мы снова минимизируем \textbf{квадратичную модель} \(m_k(s)\). У квадратичной функции на \(\R^n\) градиент по \(s\) линейен; условие \(\nabla_s m_k(s)=0\) даёт
\[
  \grad f(x_k)+H(x_k)s=0
  \quad\Leftrightarrow\quad
  H(x_k)s=-\grad f(x_k).
\]
Если \(H(x_k)\) обратима, то формально \(s_k=-H(x_k)^{-1}\grad f(x_k)\) --- это прямое обобщение формулы \(t=-f'/f''\): вместо одной второй производной стоит \textbf{обратная матрица} кривизн.

Обновление точки:
\[
  x_{k+1}=x_k+s_k.
\]

\subsubsection{Как это считают на практике (и почему не ``инвертируют Гессиан'')}
Обычно \textbf{не вычисляют} \(H^{-1}\) явно: это дорого. Решают систему линейных уравнений
\[
  H(x_k)s_k=-\grad f(x_k)
\]
методом Холецкого/LU/итерационно --- получают только вектор \(s_k\).

\subsubsection{Когда шаг Ньютона ``хороший'', а когда опасный}
\begin{itemize}
  \item Если \(H(x_k)\succ 0\) (все собственные значения \(>0\)), то \(s_k\) --- направление \textbf{спуска} (для малых шагов \(f\) убывает вдоль \(s_k\)), и квадратичная модель действительно имеет минимум.
  \item Если \(H(x_k)\) не знакопостоянна (есть отрицательные собственные значения), ``минимум модели'' может быть седловым или не тем, что нужно: направление может увести к \textbf{максимуму} или седлу исходной \(f\).
  \item Если \(H\) близка к вырожденной, решение СЛАУ неустойчиво.
\end{itemize}
Поэтому в коде почти всегда стоят \textbf{демпфирование} (шаг \(x_{k+1}=x_k+\alpha_k s_k\) с подбором \(\alpha_k\)) и/или \textbf{сдвиг} матрицы \(H+\lambda I\), и/или доверительные области (trust region).

\textbf{Плюсы:} около решения, где модель близка к правде и \(H\succ 0\), сходимость очень быстрая (``квадратичная'' по ошибке --- на каждом шаге число верных цифр растёт примерно вдвое).

\textbf{Минусы:} цена шага \(O(n^3)\) при плотном \(H\) наивным решением СЛАУ, плюс нужно получать вторые производные (аналитически, автодифференцированием или конечными разностями).

\textit{Удобное вводное чтение (в т.\,ч. вывод \(t=-f'(x_k)/f''(x_k)\) и переход к \(\R^n\)):}
\href{https://en.wikipedia.org/wiki/Newton\%27s_method_in_optimization}{Wikipedia: Newton's method in optimization}.

\subsection{Модификации Ньютона (что именно улучшают)}
\begin{itemize}
  \item \textbf{Дампированный Ньютон:} \(x_{k+1}=x_k+\alpha_k s_k\), где \(\alpha_k\) подбирают по Армихо/Вольфе.
  \item \textbf{Сдвиг (Levenberg--Marquardt идея):} \((H(x_k)+\lambda_k I)s_k=-\grad f(x_k)\), чтобы матрица была лучше обусловлена.
  \item \textbf{Усечённый Ньютон:} решают \(H s=-g\) приближённо итерациями, не формируя \(H\) явно (матвекторы).
\end{itemize}

\subsection{Квазиньютоновские методы (DFP, BFGS) --- полная запись}
\textbf{Зачем.} Метод Ньютона требует на каждом шаге точный Гессиан \(\nabla^2 f(x_k)\) и решение СЛАУ. Квазиньютоновские методы хранят \textbf{приближение} кривизны и обновляют его дёшево (поправка ранга 1--2), используя только градиенты в соседних точках.

\textbf{Обозначения.}
\begin{itemize}
  \item \(s_k := x_{k+1}-x_k\) --- фактический сдвиг на шаге;
  \item \(y_k := \grad f(x_{k+1})-\grad f(x_k)\) --- изменение градиента;
  \item \(B_k\) --- симметричная аппроксимация \textbf{Гессиана}; \(H_k:=B_k^{-1}\) --- аппроксимация \textbf{обратного} Гессиана (в коде BFGS часто обновляют именно \(H_k\)).
\end{itemize}

\textbf{Условие секущей.} Требуют, чтобы новая матрица согласовалась с наблюдением ``\(y\approx B s\)'' на последнем шаге:
\[
  B_{k+1}s_k=y_k.
\]
Для квадратичной \(f\) с матрицей \(Q\) равенство \(y=Qs\) выполнялось бы точно; для общей функции мы подгоняем только последний шаг.

\textbf{Почему нужно \(y_k^\top s_k>0\).} Для выпуклой квадратичной функции это выполняется строго. При линейном поиске по правилам \textbf{Вольфе} для минимизации это неравенство обычно сохраняется --- тогда обновления DFP/BFGS сохраняют положительную определённость \(B_k\) (при \(B_0\succ0\)).

\subsubsection{Каркас алгоритма квазиньютоновского спуска}
\begin{enumerate}
  \item Задать \(x_0\), симметричную \(B_0\succ0\) (часто \(B_0=I\)).
  \item Решить \(B_k p_k=-\grad f(x_k)\) (или \(p_k=-H_k\grad f(x_k)\), если храните \(H_k\)).
  \item Линейный поиск: найти \(\alpha_k\), положить \(s_k=\alpha_k p_k\), \(x_{k+1}=x_k+s_k\).
  \item \(y_k=\grad f(x_{k+1})-\grad f(x_k)\). Если \(y_k^\top s_k\le0\), не обновлять матрицу или уменьшать шаг.
  \item Обновить \(B_k\) или \(H_k\) по DFP/BFGS.
\end{enumerate}

\subsubsection{DFP (обновление аппроксимации Гессиана \(B_k\))}
Пусть \(y_k^\top s_k>0\) и \(s_k^\top B_k s_k>0\). Стандартная форма:
\[
  B_{k+1}=B_k+\frac{y_k y_k^\top}{y_k^\top s_k}-\frac{B_k s_k s_k^\top B_k}{s_k^\top B_k s_k}.
\]

\subsubsection{BFGS (обновление \(B_k\))}
\[
  B_{k+1}=B_k-\frac{B_k s_k s_k^\top B_k}{s_k^\top B_k s_k}+\frac{y_k y_k^\top}{y_k^\top s_k}.
\]

\subsubsection{BFGS (обновление \(H_k\approx B_k^{-1}\), удобно в реализации)}
Положим \(\rho_k:=1/(y_k^\top s_k)\). Тогда
\[
  H_{k+1}=\bigl(I-\rho_k s_k y_k^\top\bigr)\,H_k\,\bigl(I-\rho_k y_k s_k^\top\bigr)+\rho_k\,s_k s_k^\top.
\]
Направление \(p_{k+1}=-H_{k+1}\grad f(x_{k+1})\) получается без решения СЛАУ (стоимость \(O(n^2)\) при плотной \(H_k\)).

\subsubsection{Сложность и L-BFGS}
Плотное хранение \(B_k\) или \(H_k\): память \(O(n^2)\). Для больших \(n\) используют \textbf{L-BFGS}: хранят только \(m\ll n\) последних пар \((s_j,y_j)\) и восстанавливают действие \(H_k g\) за \(O(mn)\).

\textit{Справочно:} \href{https://en.wikipedia.org/wiki/Broyden-Fletcher-Goldfarb-Shanno_algorithm}{Wikipedia: BFGS}.

%--------------------------------------------------------------------
\section[Без производных]{Методы без производных и покоординатный спуск}
%--------------------------------------------------------------------

\subsection{Метод Гаусса--Зейделя (покоординатный спуск) --- как читать формулу}
\textbf{Зачем.} Свести сложную задачу в \(\R^n\) к серии простых одномерных минимизаций.

\textbf{В чём путаница.} В короткой записи
\[
  x^{(i)} \leftarrow \arg\min_{\alpha} f\bigl(x^{(1)},\ldots,x^{(i-1)},\alpha,x^{(i+1)},\ldots,x^{(n)}\bigr)
\]
не видно, какие \(x^{(j)}\) уже \textbf{новые}, а какие ещё \textbf{старые}. В методе \textbf{Гаусса--Зейделя} при минимизации по \(i\)-й координате используют уже обновлённые координаты \(1,\ldots,i-1\) и ещё не тронутые \(i+1,\ldots,n\).

\textbf{Ясная запись через промежуточные векторы \(z^{(i)}\).}
Пусть \(x^{[t]}\) --- вектор перед очередным полным проходом. Положим \(z^{(0)}=x^{[t]}\). Для каждого \(i=1,\ldots,n\) определим \(z^{(i)}\in\R^n\) так:
\[
  z^{(i)}_j=
  \begin{cases}
    z^{(i-1)}_j, & j\neq i,\\[2mm]
    \arg\min\limits_{\alpha\in\R} f\bigl(z^{(i-1)}_1,\ldots,z^{(i-1)}_{i-1},\alpha,z^{(i-1)}_{i+1},\ldots,z^{(i-1)}_{n}\bigr), & j=i.
  \end{cases}
\]
После цикла \(i=1,\ldots,n\) кладут \(x^{[t+1]}:=z^{(n)}\).

\textbf{Отличие от метода Якоби.} В методе Якоби все координаты на одном проходе считают из \emph{одного} \(x^{[t]}\), затем разом делают \(x^{[t+1]}\). У Гаусса--Зейделя информация ``протекает'' по координатам сразу --- часто быстрее, но сложнее распараллелить.

\textbf{Остановка.} \(\norm{x^{[t+1]}-x^{[t]}}\le\varepsilon\) или малость \(\bigl|f(x^{[t+1]})-f(x^{[t]})\bigr|\).

\textbf{Когда хорош.} Частично разделимые \(f\), дешёвые одномерные подзадачи. \textbf{Когда плох.} Сильные ``овраги'' --- как у градиентного спуска, много проходов без большого эффекта.

\subsection{Метод Хука--Дживса (pattern search)}
\textbf{Зачем.} Минимизировать \(f:\R^n\to\R\) без градиента; допустимы шум и ``чёрный ящик''. Хорош при \(n\) небольшом (десятки).

\textbf{Идея.} Две фазы чередуются:
\begin{itemize}
  \item \textbf{Разведка (exploratory move):} из текущей базовой точки пробуют маленькие смещения \(\pm\Delta\) по каждой оси по очереди. Если по оси \(i\) нашлось улучшение, \(i\)-я координата остаётся лучшей; если нет --- возвращают старое значение.
  \item \textbf{Шаг по образцу (pattern move):} если разведка улучшила точку, строят направление ``куда мы вообще ушли за разведку'' \(d=x^{\mathrm{new}}-x^{\mathrm{base}}\) и пробуют точку
  \[
    x^{\mathrm{pat}}=x^{\mathrm{new}}+\gamma\,d,\qquad \gamma>0\ \text{(часто }\gamma=1\text{)}.
  \]
  Вокруг \(x^{\mathrm{pat}}\) снова делают разведку. Идея: если долина вытянута, суммарный сдвиг по осям указывает вдоль долины, и шаг по образцу ускоряет движение.
\end{itemize}

\textbf{Параметры.} Стартовый \(\Delta\), минимальный \(\Delta_{\min}\), коэффициент \(\gamma\), правило деления \(\Delta\) (часто \(\Delta\leftarrow\Delta/2\)), если разведка из базовой точки \emph{не} дала улучшения.

\textbf{Псевдокод одной ``большой'' итерации.}
\begin{enumerate}
  \item Задать базовую точку \(x_b\).
  \item \textbf{Разведка:} \(x\leftarrow x_b\); для \(i=1,\ldots,n\): если \(f(x-\Delta e_i)<f(x)\), то \(x\leftarrow x-\Delta e_i\); иначе если \(f(x+\Delta e_i)<f(x)\), то \(x\leftarrow x+\Delta e_i\).
  \item Если \(f(x)<f(x_b)\): положить \(d=x-x_b\), \(x_{\mathrm{trial}}=x+\gamma d\); выполнить разведку от \(x_{\mathrm{trial}}\); если улучшение есть, принять новую базу, иначе базой остаётся \(x\).
  \item Если на шаге 2 не было улучшения относительно \(x_b\), уменьшить \(\Delta\).
\end{enumerate}

\textbf{Остановка.} \(\Delta\le\Delta_{\min}\) или исчерпан бюджет вызовов \(f\).

\textbf{Плюсы/минусы.} Плюс: простота, нет производных. Минус: при больших \(n\) и оврагах может быть очень медленным; нет глобальной сходимости в общем случае.

\subsection{Метод Пауэлла (метод сопряжённых направлений без градиента)}
\textbf{Зачем.} Улучшить покоординатный спуск: в ``овраге'' последовательные минимизации по осям дают зигзаги; Пауэлл добавляет направление ``куда мы реально продвинулись за весь цикл''.

\textbf{Данные.} Начальная точка \(x_0\), линейно независимые направления \(d_1,\ldots,d_n\) (обычно \(d_i=e_i\)).

\textbf{Один цикл (основная схема).}
\begin{enumerate}
  \item Положить \(z_0:=x\) (текущая лучшая точка).
  \item Для \(i=1,\ldots,n\): найти \(\alpha_i\in\R\), минимизирующее \(\varphi_i(\alpha)=f(z_{i-1}+\alpha d_i)\) (одномерный поиск), и положить
  \[
    z_i:=z_{i-1}+\alpha_i d_i.
  \]
  \item Суммарный сдвиг за цикл:
  \[
    s:=z_n-z_0=\sum_{i=1}^n \alpha_i d_i.
  \]
  \item \textbf{Обновление направлений (классическая идея Пауэлла).} Удалить одно из старых направлений (часто то, вдоль которого было \emph{наименьшее} уменьшение \(f\) --- интуиция: оно ``менее полезно''), а на его место поставить \(s\) (после нормализации или без --- в разных изложениях по-разному). Остальные направления сдвинуть по порядку, чтобы снова было ровно \(n\) направлений.
  \item Положить \(x\leftarrow z_n\) и перейти к следующему циклу.
\end{enumerate}

\textbf{Важное замечание.} В оригинальном методе Пауэлла нужна проверка, чтобы направления не стали почти линейно зависимыми (иначе одномерные поиски ``ломаются''). В учебных конспектах это часто опускают, но в реализации добавляют критерий или перезапускают направления на оси.

\textbf{Остановка.} \(\norm{s}\le\varepsilon\) или малость прироста \(f\).

\subsection{Нелдер--Мид: симплекс в \(\R^n\)}
\textbf{Зачем.} То же, что Хук--Дживс: нет градиента, \(f\) --- чёрный ящик; размерность обычно небольшая.

\textbf{Геометрия.} В \(\R^n\) берут \(n+1\) точку \(x_1,\ldots,x_{n+1}\) в общем положении (не лежащие в одной гиперплоскости). Их выпуклая оболочка --- симплекс (сегмент при \(n=1\), треугольник при \(n=2\), \ldots).

\textbf{Обозначения на итерации.} Пусть вершины пронумерованы так, что
\[
  f(x_1)\le f(x_2)\le\cdots\le f(x_{n+1}),
\]
то есть \(x_{n+1}\) --- \textbf{худшая}, \(x_1\) --- \textbf{лучшая}. Центр ``хорошей'' грани (без худшей вершины):
\[
  \bar x=\frac{1}{n}\sum_{i=1}^n x_i.
\]

\textbf{Отражение (reflection).} Параметр \(\rho>0\) (часто \(1\)):
\[
  x_r=\bar x+\rho(\bar x-x_{n+1}).
\]
Если \(f(x_r)\) лучше, чем у второй по плохости, но не лучше \(f(x_1)\), заменяют \(x_{n+1}\leftarrow x_r\).

\textbf{Растяжение (expansion).} Если \(f(x_r)<f(x_1)\) (отражение дало новый рекорд), пробуют
\[
  x_e=\bar x+\chi(x_r-\bar x),\qquad \chi>1\ \text{(часто }2\text{)}.
\]
Если \(f(x_e)\) ещё лучше, \(x_{n+1}\leftarrow x_e\), иначе \(x_{n+1}\leftarrow x_r\).

\textbf{Сжатие внешнее (outside contraction).} Если \(x_r\) не лучше второй худшей, но лучше старой худшей, иногда делают
\[
  x_c=\bar x+\gamma(x_r-\bar x),\qquad 0<\gamma<1.
\]

\textbf{Сжатие внутреннее (inside contraction).} Если \(x_r\) хуже всех кроме старой худшей, сжимают ``внутрь'':
\[
  x_{cc}=\bar x-\gamma(\bar x-x_{n+1}).
\]

\textbf{Глобальное уменьшение (shrink).} Если сжатия не помогли, все вершины кроме лучшей стягивают к \(x_1\):
\[
  x_i\leftarrow x_1+\sigma(x_i-x_1),\quad i=2,\ldots,n+1,\quad 0<\sigma<1\ \text{(часто }0{,}5\text{)}.
\]

\textbf{Остановка.} Малость симплекса \(\max_i \norm{x_i-x_1}\le\varepsilon\) или малость разброса значений \(f\).

\textbf{Плюсы/минусы.} Плюс: очень простая логика, мало кода. Минус: нет общих гарантий сходимости; при большом \(n\) обычно плох; при ``плоской'' \(f\) симплекс вырождается.

%--------------------------------------------------------------------
\section[Штрафы и барьеры]{Условная оптимизация: штрафы и барьеры}
%--------------------------------------------------------------------

\subsection{Метод штрафных функций}
\textbf{Исходная задача.} \(\min_{x\in\Omega} f(x)\), где \(\Omega=\{x:\ g_i(x)\le0,\ h_j(x)=0\}\).

\textbf{Идея ``штрафа''.} Вместо жёстких ограничений вводят гладкую функцию \(P(x)\ge0\), которая равна нулю на \(\Omega\) и быстро растёт при выходе наружу. Решают последовательность задач
\[
  \min_{x\in\R^n}\ F_\rho(x)=f(x)+\rho\,P(x),\qquad \rho>0.
\]
При \(\rho\to+\infty\) минимизаторы \(F_\rho\) обычно стремятся к решению исходной задачи (при стандартных предположениях).

\textbf{Пример квадратичного штрафа для неравенств и равенств.}
\[
  P(x)=\sum_{i=1}^m \bigl(\max\{0,g_i(x)\}\bigr)^2+\sum_{j=1}^p \bigl(h_j(x)\bigr)^2.
\]
Тогда вне \(\Omega\) добавка \(\rho P\) ``штрафует'' за нарушение: чем дальше за границу, тем больше \(P\).

\textbf{Градиент (если нужен для BFGS).} Для одного слагаемого \(\bigl(\max\{0,g\}\bigr)^2\) в точках, где \(g(x)<0\), производная ноль; где \(g(x)>0\), \(\nabla=2g(x)\grad g(x)\). На границе \(g=0\) гладкость теряется --- на практике работают почти всегда вне ``осколков''.

\textbf{Внешний цикл (как программируют).}
\begin{enumerate}
  \item Выбрать \(\rho_1\) (не слишком большое), старт \(x^{(0)}\).
  \item Для \(s=1,2,\ldots\): решить \(\min_x F_{\rho_s}(x)\) приближённо, взяв старт с решения при \(\rho_{s-1}\); увеличить \(\rho_{s+1}=\sigma\rho_s\) (\(\sigma>1\)).
  \item Остановиться, когда \(\max_i g_i(x)\le\varepsilon_g\) и \(|h_j(x)|\le\varepsilon_h\).
\end{enumerate}

\textbf{Минусы.} При больших \(\rho\) задача становится \textbf{плохо обусловленной} (узкие овраги вдоль границы), градиенты огромны --- нужны устойчивые оптимизаторы и иногда регуляризация.

\subsection{Метод барьерных функций (внутренние точки)}
\textbf{Идея.} Для строгих неравенств \(g_i(x)<0\) строят вспомогательную задачу \emph{только для внутренних} точек:
\[
  \min_{x\in\mathrm{int}\,\Omega}\ B_\mu(x),\qquad B_\mu(x):=f(x)-\mu\sum_{i=1}^m \ln\bigl(-g_i(x)\bigr),\quad \mu>0.
\]
Логарифм \(\ln(-g_i)\) стремится к \(+\infty\), когда \(g_i\to0^{-}\) --- как ``невидимая стена'' у границы.

\textbf{Поведение при \(\mu\to0\).} Чем меньше \(\mu\), тем слабее барьер далеко от границы, и минимум \(B_\mu\) ближе к решению исходной задачи (часто к точке на границе).

\textbf{Градиент барьерной задачи.}
\[
  \grad B_\mu(x)=\grad f(x)-\mu\sum_{i=1}^m \frac{\grad g_i(x)}{g_i(x)}.
\]
У границы члены \(\mu/|g_i|\) велики --- шаги делают малыми.

\textbf{Связь с внутренне-точечными методами.} Логарифмический барьер --- ключевая идея primal-dual interior-point методов для ЛП и выпуклой оптимизации.

%--------------------------------------------------------------------
\section{Линейное программирование (ЛП)}
%--------------------------------------------------------------------
\textbf{Задача.} Линейная целевая функция и линейные ограничения; стандартная форма для интуиции:
\[
  \min_{x\in\R^n} c^\top x\quad\text{при}\quad Ax\le b,\ x\ge 0.
\]
Здесь \(c\in\R^n\), \(A\in\R^{m\times n}\), \(b\in\R^m\).

\textbf{Геометрия.} Множество \(\{x:Ax\le b,\ x\ge0\}\) --- выпуклый \textbf{многогранник} (пересечение конечного числа полупространств). Линейная функция на многограннике достигает экстремума в вершине (если многогранник ограничен в направлении минимизации).

\textbf{Симплекс-метод (словами).} Стартуют из допустимой вершины. На каждом шаге выбирают \textbf{входящую} переменную (по отрицательной ``оценке'' в строке цели) и \textbf{выходящую} (минимальное положительное отношение правых частей к положительным коэффициентам --- правило минимального отношения), чтобы перейти по ребру в соседнюю вершину с не большей целевой функцией. Пока есть улучшение --- продолжают.

\textbf{Табличная форма.} На зачёте часто просят уметь читать симплекс-таблицу: базисные переменные, приведённая целевая строка, допустимость.

\textbf{Внутренне-точечные методы.} Идут через \emph{внутренность} допустимого множества (барьеры/примал--дуальные шаги), часто полиномиальная сложность в теории; на практике конкурируют с симплексом на больших разреженных задачах.

\textit{Справочно:} \href{https://en.wikipedia.org/wiki/Simplex_algorithm}{Wikipedia: Simplex algorithm}.

%--------------------------------------------------------------------
\section[Глобальная оптимизация]{Глобальная оптимизация: PSO и генетические алгоритмы}
%--------------------------------------------------------------------
Это \textbf{эвристики}: глобальный минимум не гарантирован, но на практике часто находят хорошее решение при нескольких запусках и настройке параметров.

\subsection{Рой частиц (PSO, Particle Swarm Optimization)}
\textbf{Биологическая метафора.} Стая ``летит'' по пространству параметров; каждая частица помнит свой лучший найденный пункт, а также знает (или знает у соседей) лучшую точку роя.

\textbf{Состояние.} Частица с номером \(i\): позиция \(x_i\in\R^n\), скорость \(v_i\in\R^n\). Индивидуальный лучший минимум \(p_i\) (personal best), глобальный лучший \(g\) (global best). В варианте \textbf{lbest} вместо \(g\) используют лучшую точку среди соседей по заданной топологии.

\textbf{Классическое обновление (инерционная модель).} На шаге \(t\):
\begin{align*}
  v_i^{t+1} &= \omega v_i^t + c_1 r_1^t\odot(p_i-x_i^t)+c_2 r_2^t\odot(g-x_i^t),\\
  x_i^{t+1} &= x_i^t+v_i^{t+1},
\end{align*}
где \(r_1^t,r_2^t\) --- векторы из \(U(0,1)\) по координатам (поэлементное умножение \(\odot\)), \(\omega,c_1,c_2>0\).

\textbf{Смысл коэффициентов.} \(\omega\) --- инерция (больше \(\omega\) --- больше ``разлёт'' и исследование). \(c_1\) --- вес ``к своей памяти''; \(c_2\) --- вес ``к общему лидеру''.

\textbf{Ограничения.} Часто вводят \(v_i\leftarrow \mathrm{clip}(v_i,-v_{\max},v_{\max})\) и \(x_i\) режут по допустимой области (проекция или штраф).

\textbf{Псевдокод.}
\begin{enumerate}
  \item Инициализировать \(x_i,v_i\), вычислить \(f(x_i)\), задать \(p_i=x_i\), найти \(g\).
  \item Повторять: обновить \(v_i,x_i\); если \(f(x_i)<f(p_i)\), то \(p_i\leftarrow x_i\); если \(f(x_i)<f(g)\), то \(g\leftarrow x_i\); критерий остановки.
\end{enumerate}

\subsection{Генетический алгоритм (ГА)}
\textbf{Кодирование.} \textbf{Бинарное:} хромосома --- строка битов; \textbf{вещественное:} хромосома --- вектор \(x\in\R^n\).

\textbf{Fitness.} Для минимизации \(f\) задают, например, \(\mathrm{fitness}(x)=-f(x)\) (чем больше fitness, тем лучше).

\textbf{Турнирная селекция размера \(T\).} Для выбора одного родителя: взять \(T\) случайных особей, выбрать лучшую по fitness. Повторить, чтобы набрать пару родителей.

\textbf{Кроссинговер (примеры).}
\begin{itemize}
  \item \textbf{Одноточечный (бинарный):} разрезать две строки в позиции \(k\) и обменять хвосты.
  \item \textbf{Арифметический (вещественный):} для родителей \(a,b\in\R^n\) и случайного \(\lambda\in[0,1]\) потомок \(c=\lambda a+(1-\lambda)b\) (иногда два потомка симметрично).
\end{itemize}

\textbf{Мутация.} Бинарная: инвертировать каждый бит с вероятностью \(p_m\). Вещественная: \(x\leftarrow x+\sigma\xi\), \(\xi\sim\mathcal{N}(0,I)\).

\textbf{Элитизм.} Сохранить \(E\) лучших особей из поколения \(t\) в поколение \(t+1\) без изменений --- чтобы не потерять лучшее решение.

\textbf{Схема поколения.}
\begin{enumerate}
  \item Инициализация популяции; оценка fitness.
  \item Пока не стоп: селекция \(\to\) кроссинговер \(\to\) мутация \(\to\) формирование нового поколения с элитизмом.
\end{enumerate}

\textbf{Баланс.} Сильная селекция и слабая мутация --- быстрая сходимость к локальному оптимуму; слабая селекция и сильная мутация --- долгое блуждание.

%--------------------------------------------------------------------
\section[Имитация отжига]{Имитация отжига (SA): от порога к вероятности}
%--------------------------------------------------------------------
\textbf{Задача.} \(\min_{X\in\Omega}\varphi(X)\), где \(\Omega\) --- дискретное или непрерывное множество допустимых решений (перестановки, целочисленные векторы, вещественный ящик и т.д.).

\textbf{Почему не жадный локальный поиск.} Если на каждом шаге брать только улучшения \(\varphi(X')<\varphi(X)\), легко застрять в локальном минимуме. Отжиг разрешает иногда \textbf{ухудшения}, чтобы ``выпрыгнуть'' из ямы.

\textbf{Окрестность.} Для каждого \(X\in\Omega\) задано конечное (или измеримое) множество соседей \(d(X)\subseteq\Omega\). На шаге из \(X_t\) выбирают случайного соседа \(X'_t\in d(X_t)\) (равномерно или по заданным весам).

\subsection{Пороговый алгоритм (мотивация)}
Детерминированная эвристика: принимать \(X'\), если \(\varphi(X')-\varphi(X)<\varepsilon(t)\), где \(\varepsilon(t)\to0\). Тогда в начале допускаются большие ухудшения, в конце --- почти жадность. Недостаток: неясно, как выбрать \(\varepsilon(t)\) универсально.

\subsection{Правило Метрополиса (стандарт SA для минимизации)}
Положим \(\Delta_t=\varphi(X'_t)-\varphi(X_t)\) --- приращение цели (хотим минимум, значит \(\Delta_t<0\) --- улучшение). Вероятность принятия кандидата:
\[
  p_{\mathrm{acc}}(\Delta_t,T_t)=
  \begin{cases}
    1, & \Delta_t\le 0,\\[2mm]
    \exp(-\Delta_t/T_t), & \Delta_t>0,
  \end{cases}
\]
где \(T_t>0\) --- \textbf{температура}. Генерируют \(u_t\sim U(0,1)\); если \(u_t<p_{\mathrm{acc}}\), то \(X_{t+1}=X'_t\), иначе \(X_{t+1}=X_t\).

\textbf{Связь с Больцманом (для интуиции).} В статфизике вероятности состояний пропорциональны \(\exp(-E/kT)\); здесь \(\varphi\) играет роль энергии, \(T\) --- ``шум''.

\textbf{Поведение при разных \(T\).} При большом \(T_t\) даже \(\Delta_t\gg0\) даёт \(\exp(-\Delta/T)\approx 1\) --- почти всё принимается (исследование). При малом \(T_t\) ухудшения почти не проходят --- локальная доводка.

\subsection{Полный псевдокод (с лучшим найденным решением)}
\begin{enumerate}
  \item Выбрать \(X\in\Omega\), вычислить \(\varphi(X)\); положить \(X_{\mathrm{best}}\leftarrow X\), \(\varphi_{\mathrm{best}}\leftarrow\varphi(X)\).
  \item Задать \(T\leftarrow T_0\), функцию охлаждения \(T\leftarrow g(T,t)\), параметры остановки.
  \item Для \(t=1,2,\ldots\): сгенерировать \(X'\in d(X)\), \(\Delta\leftarrow\varphi(X')-\varphi(X)\); принять по правилу Метрополиса; если \(\varphi(X)<\varphi_{\mathrm{best}}\), обновить рекорд; обновить \(T\); проверить стоп.
\end{enumerate}

\subsection{Законы охлаждения (как в методичке)}
\begin{itemize}
  \item \textbf{Геометрическое:} \(T_{t}=\alpha T_{t-1}\), \(\alpha\in[0{,}8;0{,}99]\) --- простой и частый выбор на практике.
  \item \textbf{Логарифмическое:} \(T_t=T_0/\ln(t+1)\) --- теоретически связано с медленным охлаждением (на практике может быть очень медленным).
  \item \textbf{Very fast annealing:} \(T_t=T_0\exp(-\beta t^{1/|X|})\), где \(|X|\) трактуют как размерность решения.
\end{itemize}

\subsection{Практические советы}
\begin{itemize}
  \item \textbf{Окрестность} должна быть достаточно ``богатой'': из любого разумного состояния за несколько шагов можно добраться до области глобально хороших решений (на графе состояний --- связность).
  \item \textbf{Начальная температура} часто подбирают так, чтобы в начале принималась заметная доля ухудшающих шагов (например, 20--60\% для умеренных \(\Delta\)).
  \item \textbf{Остановка:} \(T_t<T_{\min}\), лимит итераций, либо нет улучшения рекорда \(K\) шагов подряд.
\end{itemize}

%--------------------------------------------------------------------
\section{Краткая памятка: что выбрать}
%--------------------------------------------------------------------
\begin{itemize}
  \item \textbf{Одна переменная, унимодальность, отрезок известен:} золотое сечение или Фибоначчи.
  \item \textbf{Отрезок неизвестен:} Свенн + золотое сечение.
  \item \textbf{Гладкая \(f\), есть \(\grad f\):} градиентный/наискорейший; для скорости и устойчивости --- BFGS с Вольфе.
  \item \textbf{Квадратичная положительно определённая \(A\):} сопряжённые градиенты.
  \item \textbf{Нет производных, мало размерность:} Нелдер--Мид или Хук--Дживс.
  \item \textbf{Ограничения ``в лоб'':} штрафы/барьеры + внутренний метод.
  \item \textbf{Много локальных минимумов / комбинаторика:} SA, ГА, PSO; часто + локальная доводка.
\end{itemize}

\end{document}
